% FILE: 06_contribution_to_thesis.tex
% PROJECT: ma — M&A Claim Taxonomy and Board Projection Surface

\section{Contribution to the Dissertation}
\label{sec:ma-contribution}

\subsection{Contribution to Prediction vs.~Coordination}
\label{sec:ma-contrib-prediction}

The central dissertation argument is that prediction---local text imitation, statistical
forecasting, AI-generated content---does not conserve consequence, and that the enterprise
requires coordination: shared, replayable, cryptographically attested process truth. The
\texttt{ma/} module provides the most economically explicit instantiation of this argument
in the thesis.

The traditional M\&A synergy model is a prediction artifact. A financial model that projects
\$50M in post-merger EBITDA improvement is a prediction in the precise sense that fails: it
asserts an organizational future without being grounded in a replayable record of the
organizational past. The prediction may be fluent---well-structured, numerically consistent,
defensible in narrative---while being entirely disconnected from the actual process execution
history of either party. When this prediction fails post-close, the buyer has paid a merger
premium for a narration.

The \texttt{ma/} framework replaces this with coordination: every claim is grounded in an
OCEL 2.0 event log that records actual execution history, verified by a conformance query
that any party can replay, and sealed by a receipt that any auditor can independently
re-execute. The EBITDA claim of \$1,250,000 in \texttt{rec\_ebitda\_rework\_001.json} is not
a prediction---it is a coordinate: a point in the space of verifiable process reality, with
fitness 0.982 and precision 0.945 certifying that the underlying event log is a faithful
record of the claimed process.

This is the dissertation's core claim demonstrated at the capital-market surface: coordination
(receipted, replayable process evidence) has a higher truth value than prediction (spreadsheet
extrapolation), and this higher truth value is expressible as a merger premium.

\subsection{Contribution to Process Evidence}
\label{sec:ma-contrib-process-evidence}

The \texttt{ma/} module makes two specific contributions to the process evidence strand of
the thesis.

\textbf{First}, it defines the evidentiary escalation ladder (T1--T6 Board Claim Taxonomy)
that translates process mining output into financial claim tiers. This ladder is a bridge
between the academic process mining literature (Van der Aalst, Adriansyah, Leemans, Weidlich,
Ghahfarokhi) and the capital-market practice of M\&A diligence. T3 Conformance claims require
the same alignment-based fitness computation that academic process mining uses for evaluation.
T4 Performance claims require event timestamp analysis at the same rigor as throughput-time
benchmarks in the academic literature. The \texttt{ma/} module operationalizes academic
process mining metrics as capital-market admissible evidence.

\textbf{Second}, it defines the Reverse Porter Five Forces inversion---the claim that process
evidence permanently restructures competitive dynamics rather than merely defending against
competitive threats. This is an original thesis contribution (AUTHOR THESIS): the argument
that a multi-year OCEL event log history is structurally irreproducible by a new entrant, and
that this irreproducibility constitutes a persistent competitive moat, is not derived from a
single prior source. It is grounded in the structural logic of retrospective data (a competitor
who starts logging today cannot produce three years of historical OCEL logs) but the
capital-market framing of this as a Porter inversion is this thesis's contribution.

\subsection{Contribution to Manufacturing Architecture}
\label{sec:ma-contrib-manufacturing}

The \texttt{ma/} module defines the board-projection surface of the CodeManufactory
architecture. The Board Projection Renderer is a seven-step manufacturing pipeline: Claims
JSON is the raw input; sealed PowerPoint and Excel artifacts are the manufactured outputs;
the intermediate stages (Tera2 templating, pptx-rs assembly, calamine workbook assembly,
receipt validation, ledger sealing) are manufacturing operators in the CodeManufactory sense.

The receipt chain is a manufacturing lineage record: every manufactured claim artifact traces
back through the receipt to the event log from which it was derived. This is the
manufacturing architecture principle---that the provenance of every artifact must be
independently verifiable from its production inputs---applied to financial claims at the
capital-market surface.

The VDR structure (\texttt{/event-logs/}, \texttt{/models/}, \texttt{/receipts/},
\texttt{/conformance/}, \texttt{/provenance/}, \texttt{/synergy/}) is a manufacturing
artifact repository: a structured archive of all intermediate and final artifacts in the
pipeline, organized for independent inspection rather than for internal convenience.

\subsection{Contribution to Capital-Flow Infrastructure}
\label{sec:ma-contrib-capital-flow}

The \texttt{ma/} module is the capital-flow layer of the thesis architecture. It is the point
at which process intelligence---which in \texttt{wasm4pm} and \texttt{ggen} operates at the
operational and ontological levels---is translated into capital-market claims, merger
premiums, and acquisition valuations.

The contribution to capital-flow infrastructure is structural. The module establishes that:

\begin{enumerate}
  \item Operational debt (process spaghetti, compliance deficits, legacy lock-in, shadow IT)
    is quantifiable as a valuation haircut: 5--30\% depending on debt class and severity. This
    makes process evidence directly relevant to the term sheet.
  \item Synergy claims are verifiable via bisimulation equivalence (Process Equivalence Model)
    and Execution Path Intersection---making the merger premium defensible to a skeptical
    buyer's auditor rather than merely claimed by the seller's investment banker.
  \item The diligence compression from 7+ weeks to 3 days, enabled by wasm4pm automated
    query execution, is a transaction-infrastructure improvement: it reduces the cost of
    capital by reducing due-diligence friction, making process-intelligent targets more
    liquid acquisition candidates.
  \item The Hostile Takeover Receipt (v30.1.1) extends this infrastructure to adversarial
    contexts, where incumbent management actively conceals operational debt. The HTR taxonomy
    establishes that ghost dependencies and exception-rate variance can be surfaced by
    replaying event streams---making process evidence a tool for hostile acquirers as well
    as friendly ones.
\end{enumerate}

These contributions position the \texttt{ma/} module as the thesis's primary demonstration
that process intelligence is not merely an operations analytics discipline but a
capital-allocation discipline---that the full-lifecycle process intelligence stack terminates
not in an operations dashboard but in a cryptographically sealed board presentation whose
every assertion is independently replayable.
